% ---------------------------------------------------------------
% TEMPLATE TCC2 - SENAC SANTO AMARO
% ---------------------------------------------------------------
% Trabalho de Conclusão de Curso 2
% Elaborado para apresentação do trabalho final do TCC
% Estrutura: Introdução, Revisão Bibliográfica, Metodologia,
%            Desenvolvimento, Resultados e Análise, Conclusão
% ---------------------------------------------------------------

\documentclass[
    12pt,                % tamanho da fonte
    openright,           % capítulos começam em página ímpar
    oneside,             % para impressão em apenas um lado do papel
    a4paper,             % tamanho do papel
    english,             % idioma adicional
    brazil               % idioma principal
]{abntex2}

% ---------------------------------------------------------------
% PACOTES
% ---------------------------------------------------------------
\usepackage{lmodern}                % Fonte Latin Modern
\usepackage[T1]{fontenc}            % Seleção de códigos de fonte
\usepackage[utf8]{inputenc}         % Codificação do documento
\usepackage{lastpage}               % Usado pela Ficha catalográfica
\usepackage{indentfirst}            % Indenta o primeiro parágrafo
\usepackage{color}                  % Controle de cores
\usepackage{graphicx}               % Inclusão de gráficos
\usepackage{microtype}              % Melhorias de justificação
\usepackage{float}                  % Para posicionamento de figuras
\usepackage{booktabs}               % Tabelas profissionais
\usepackage{multirow}               % Células mescladas em tabelas
\usepackage{longtable}              % Tabelas longas
\usepackage{amsmath,amssymb}        % Símbolos matemáticos
\usepackage{pgfgantt}               % Diagramas de Gantt
\usepackage[brazilian,hyperpageref]{backref}  % Páginas com citações
\usepackage[alf]{abntex2cite}       % Citações ABNT

% ---------------------------------------------------------------
% CONFIGURAÇÕES DO PDF
% ---------------------------------------------------------------
\makeatletter
\hypersetup{
    pdftitle={\@title},
    pdfauthor={\@author},
    pdfsubject={TCC2 - SENAC Santo Amaro},
    pdfcreator={LaTeX with abnTeX2},
    pdfkeywords={tcc2}{senac}{santo amaro}{trabalho de conclusão},
    colorlinks=true,        % false: links em caixas; true: links coloridos
    linkcolor=blue,         % cor dos links internos
    citecolor=blue,         % cor dos links para bibliografia
    filecolor=magenta,      % cor dos links para arquivos
    urlcolor=blue,
    bookmarksdepth=4
}
\makeatother

% ---------------------------------------------------------------
% CONFIGURAÇÕES DE APARÊNCIA DO PDF FINAL
% ---------------------------------------------------------------
\definecolor{blue}{RGB}{41,5,195}

% ---------------------------------------------------------------
% INFORMAÇÕES DO DOCUMENTO
% ---------------------------------------------------------------
\titulo{Título do Trabalho de Conclusão de Curso}
\autor{Nome Completo do Aluno}
\local{São Paulo}
\data{2025}
\orientador{Prof. Nome do Orientador}
% \coorientador{Prof. Nome do Coorientador} % Descomente se houver coorientador

\instituicao{%
  SENAC -- Serviço Nacional de Aprendizagem Comercial
  \par
  Unidade Santo Amaro
  \par
  Curso de [Nome do Curso]
}

\tipotrabalho{Trabalho de Conclusão de Curso II}

\preambulo{Trabalho de Conclusão de Curso II apresentado ao SENAC Santo Amaro como requisito parcial para conclusão do curso de [Nome do Curso].}

% ---------------------------------------------------------------
% COMPILA O ÍNDICE
% ---------------------------------------------------------------
\makeindex

% ---------------------------------------------------------------
% INÍCIO DO DOCUMENTO
% ---------------------------------------------------------------
\begin{document}

% Seleciona o idioma do documento
\selectlanguage{brazil}

% Retira espaço extra obsoleto entre as frases
\frenchspacing

% ---------------------------------------------------------------
% ELEMENTOS PRÉ-TEXTUAIS
% ---------------------------------------------------------------
\pretextual

% ---------------------------------------------------------------
% CAPA
% ---------------------------------------------------------------
\imprimircapa

% ---------------------------------------------------------------
% FOLHA DE ROSTO
% ---------------------------------------------------------------
\imprimirfolhaderosto

% ---------------------------------------------------------------
% RESUMO
% ---------------------------------------------------------------
\begin{resumo}
Este documento apresenta o Trabalho de Conclusão de Curso (TCC2) desenvolvido no SENAC Santo Amaro, com base na estrutura recomendada por Raul Sidnei Wazlawick em \textit{Metodologia de Pesquisa para Ciência da Computação}. O foco está em evidenciar o problema, os objetivos, o método, o desenvolvimento, os resultados e a contribuição científica do trabalho, evitando textos meramente descritivos. O documento está organizado em capítulos que apresentam a introdução ao tema, a revisão bibliográfica que fundamenta a pesquisa, a metodologia adotada, o desenvolvimento da solução, os resultados e a respectiva análise crítica, seguidos da conclusão e das referências bibliográficas.

\vspace{\onelineskip}

\noindent
\textbf{Palavras-chave}: palavra1. palavra2. palavra3. palavra4. palavra5.
\end{resumo}

% ---------------------------------------------------------------
% LISTA DE ILUSTRAÇÕES
% ---------------------------------------------------------------
% \pdfbookmark[0]{\listfigurename}{lof}
% \listoffigures*
% \cleardoublepage

% ---------------------------------------------------------------
% LISTA DE TABELAS
% ---------------------------------------------------------------
% \pdfbookmark[0]{\listtablename}{lot}
% \listoftables*
% \cleardoublepage

% ---------------------------------------------------------------
% LISTA DE ABREVIATURAS E SIGLAS
% ---------------------------------------------------------------
% Esta lista deve vir antes do início do trabalho (elementos pré-textuais).
% Em ordem alfabética, relacione APENAS as abreviaturas e siglas que
% realmente aparecem no texto. Edite a tabela conforme o seu trabalho.
\pdfbookmark[0]{Lista de Abreviaturas e Siglas}{loas}
\chapter*{Lista de Abreviaturas e Siglas}

\noindent A lista a seguir reúne, em ordem alfabética, as abreviaturas e siglas utilizadas neste trabalho, acompanhadas de seus respectivos significados. Substitua, acrescente ou remova linhas conforme a necessidade do seu projeto.

\vspace{\onelineskip}

\begin{center}
\begin{tabular}{@{}p{3.2cm}p{0.6\textwidth}@{}}
\toprule
\textbf{Abreviatura/Sigla} & \textbf{Significado} \\
\midrule
ABNT  & Associação Brasileira de Normas Técnicas \\
API   & \textit{Application Programming Interface} (Interface de Programação de Aplicações) \\
HTTP  & \textit{Hypertext Transfer Protocol} \\
REST  & \textit{Representational State Transfer} \\
SQL   & \textit{Structured Query Language} \\
TCC   & Trabalho de Conclusão de Curso \\
UML   & \textit{Unified Modeling Language} \\
\bottomrule
\end{tabular}
\end{center}

\cleardoublepage

% ---------------------------------------------------------------
% SUMÁRIO
% ---------------------------------------------------------------
\pdfbookmark[0]{\contentsname}{toc}
\tableofcontents*
\cleardoublepage

% ---------------------------------------------------------------
% ELEMENTOS TEXTUAIS
% ---------------------------------------------------------------
\textual

% ===============================================================
% CAPÍTULO 1 - INTRODUÇÃO
% ===============================================================
\chapter{Introdução}
\label{cap:introducao}

A introdução deve apresentar o tema do trabalho de forma clara e objetiva, situando o leitor no contexto da pesquisa. Neste capítulo são apresentados a contextualização do problema, a sua formulação, os objetivos e a hipótese que norteiam a pesquisa, bem como a justificativa para a realização do estudo.

% ---------------------------------------------------------------
\section{Contextualização do problema}
\label{sec:contexto}

Descreva o cenário atual e a lacuna identificada, embasando o texto em dados, estatísticas, políticas ou práticas profissionais. Evite narrativas históricas extensas e mantenha o foco no problema, ressaltando a relevância do assunto para a Computação e para a sociedade.

Descreva o ambiente ou domínio em que o problema ocorre, utilizando dados e referências que comprovem a sua existência. Explique quem são os principais afetados e quais são as consequências de não resolvê-lo.

\vspace{0.3cm}
\noindent\fbox{\parbox{\dimexpr\textwidth-2\fboxsep-2\fboxrule}{%
\textbf{Orientação sobre citações (ABNT NBR 10520:2023)}

\vspace{0.2cm}
O aluno deve \textbf{priorizar citações indiretas} (paráfrases), nas quais a ideia do autor é reescrita com as próprias palavras. Citações diretas (transcrições literais) \textbf{não são recomendadas}, exceto quando o texto original for insubstituível --- por exemplo, definições consagradas, leis, normas técnicas ou trechos cuja reformulação alteraria o sentido.

\vspace{0.2cm}
\textbf{Citação indireta} (recomendada) --- o autor interpreta e reescreve a ideia:

\vspace{0.1cm}
\textit{Exemplo 1 --- sistema autor-data entre parênteses:}\\
A definição clara do problema é considerada um fator determinante para o sucesso de qualquer pesquisa na área de Computação \cite{Wazlawick2020}.

\vspace{0.1cm}
\textit{Exemplo 2 --- autor como parte do texto:}\\
Para \citeonline{Silva2022}, a escolha de uma metodologia adequada impacta diretamente na qualidade dos resultados obtidos em projetos de software.

\vspace{0.2cm}
\textbf{Citação direta curta} (até 3 linhas, apenas quando indispensável) --- entre aspas, com indicação de página:

\vspace{0.1cm}
Segundo \citeonline[p.~42]{Wazlawick2020}, ``a formulação do problema deve ser expressa de forma clara, objetiva e delimitada, evitando ambiguidades que possam comprometer o escopo da pesquisa''.

\vspace{0.2cm}
\textbf{Citação direta longa} (mais de 3 linhas, apenas quando indispensável) --- recuo de 4\,cm, fonte menor, sem aspas:

\begin{citacao}
A revisão bibliográfica não consiste em simplesmente listar os trabalhos encontrados, mas em analisá-los criticamente, identificando convergências, divergências e lacunas que justifiquem a realização de uma nova pesquisa \cite[p.~87]{Wazlawick2020}.
\end{citacao}

\textbf{Resumo}: use citações indiretas como regra. Reserve a citação direta para casos excepcionais em que a transcrição literal seja realmente necessária.
}}

% ---------------------------------------------------------------
\section{Formulação do problema}
\label{sec:problema}

Coloque o problema em forma de pergunta clara, delimitando o que será investigado e qual resultado se espera atingir dentro da Computação. A pergunta de pesquisa deve ser específica o suficiente para ser respondida ao longo do trabalho e diretamente derivada da lacuna apresentada na contextualização.

% ---------------------------------------------------------------
\section{Objetivos e hipótese}
\label{sec:objetivos}

Os objetivos definem o que o trabalho pretende alcançar. Devem ser claros, mensuráveis e diretamente relacionados ao problema identificado.

\subsection{Objetivo geral}

Declare o propósito principal do TCC2 em uma frase direta, conectando-o ao problema e deixando claro se o foco é compreender, comparar, validar ou propor algo novo.

\textbf{Dica}: utilize verbos no infinitivo como \textit{desenvolver}, \textit{analisar}, \textit{propor}, \textit{avaliar}, \textit{implementar}, \textit{comparar}, entre outros. Evite verbos vagos como \textit{estudar} ou \textit{entender}.

\subsection{Objetivos específicos}

Liste objetivos mensuráveis que estruturam o plano de trabalho. Tente utilizar verbos como revisar, comparar, modelar, quantificar ou validar. Exemplos:

\begin{enumerate}
    \item Revisar obras-chave sobre [tema] e identificar lacunas;
    \item Modelar ou prototipar a abordagem proposta;
    \item Implementar experimentos controlados;
    \item Avaliar o desempenho por métricas (ex.: acurácia, latência).
\end{enumerate}

\textbf{Atenção}: cada objetivo específico deve contribuir diretamente para o alcance do objetivo geral. Evite listar objetivos que não tenham relação clara com o problema.

\subsection{Hipótese ou critério de validação}

Descreva a hipótese que será testada ou o critério usado para afirmar que o objetivo geral foi atingido. Caso não haja hipótese formal, explique como as métricas escolhidas demonstrarão o sucesso esperado.

% ---------------------------------------------------------------
\section{Justificativa}
\label{sec:justificativa}

Explicite por que resolver esse problema é importante. A justificativa deve abordar:

\begin{itemize}
    \item \textbf{Relevância acadêmica}: como este trabalho contribui para o avanço do conhecimento na área;
    \item \textbf{Relevância prática/profissional}: quais benefícios práticos a solução trará para usuários, organizações ou a sociedade;
    \item \textbf{Lacuna identificada}: o que ainda não foi resolvido ou explorado suficientemente pelos trabalhos existentes.
\end{itemize}

Ressalte a importância do problema para a Computação, destacando contribuições acadêmicas, sociais ou industriais. Explique como a proposta evita se tornar um ``manual de ferramenta''.

% ===============================================================
% CAPÍTULO 2 - REVISÃO BIBLIOGRÁFICA / TRABALHOS RELACIONADOS
% ===============================================================
\chapter{Revisão Bibliográfica / Trabalhos Relacionados}
\label{cap:revisao}

A revisão bibliográfica tem como finalidade apresentar o embasamento teórico e técnico do trabalho, demonstrando o conhecimento do autor sobre o tema e identificando o estado da arte. Este capítulo está organizado em três seções: os \textbf{fundamentos teóricos} necessários à compreensão da proposta, as \textbf{técnicas e abordagens existentes} no contexto do problema e os \textbf{trabalhos correlatos} que enfrentaram o mesmo desafio, com suas respectivas limitações.

% ---------------------------------------------------------------
\section{Fundamentos teóricos necessários}
\label{sec:fundamentos}

Explique os conceitos essenciais, definindo termos e equações quando necessário. Utilize blocos de texto ou pequenos quadros comparativos para sintetizar modelos. Organize esta seção por temas ou conceitos, não por autor, partindo do mais geral para o mais específico.

Utilize citações para fundamentar cada conceito apresentado. Uma citação direta pode ser empregada, por exemplo, com \citeonline{Silva2022} para reforçar que determinados métodos exigem validação empírica; uma citação indireta pode ser escrita como ``vários estudos destacam \ldots'' seguida de \cite{Wazlawick2020}.

% ---------------------------------------------------------------
\section{Técnicas, métodos ou abordagens existentes}
\label{sec:tecnicas}

Apresente abordagens já empregadas no contexto e destaque o que funcionou ou falhou. Inclua uma tabela comparativa simples para clarificar diferenças:

\begin{table}[htb]
  \centering
  \caption{Comparativo de abordagens relacionadas}
  \label{tab:abordagens}
  \begin{tabular}{p{4cm}p{3cm}p{4cm}}
    \toprule
    \textbf{Abordagem} & \textbf{Domínio} & \textbf{Limitação principal} \\
    \midrule
    Técnica A & IoT & Falta de validação em larga escala \\
    Técnica B & Redes neurais & Alto custo computacional \\
    \bottomrule
  \end{tabular}
  \legend{Fonte: Elaborado pelo autor (2025)}
\end{table}

% ---------------------------------------------------------------
\section{Trabalhos correlatos e limitações}
\label{sec:trabalhos_correlatos}

Liste artigos diretamente relacionados e explique por que cada um não resolve totalmente o problema (ex.: falta de comparação, escopo diferente, dados limitados). Relacione essas lacunas com o que sua pesquisa pretende preencher. Para cada trabalho correlato, descreva:

\begin{enumerate}
    \item \textbf{Qual problema} o trabalho tentou resolver (e por que é o mesmo problema que o seu);
    \item A metodologia ou abordagem utilizada para resolvê-lo;
    \item Os principais resultados obtidos;
    \item As limitações identificadas na solução proposta;
    \item Como o \textbf{seu trabalho se diferencia} ou avança em relação a ele.
\end{enumerate}

\textit{Evite a ``síndrome da intersecção esquecida'': inclua trabalhos que aplicaram a técnica na mesma área.}

% ===============================================================
% CAPÍTULO 3 - METODOLOGIA
% ===============================================================
\chapter{Metodologia}
\label{cap:metodologia}

A metodologia descreve, de forma rigorosa e reprodutível, como a pesquisa será conduzida. Este capítulo apresenta o tipo de pesquisa e a hipótese, os procedimentos da abordagem proposta, as ferramentas e tecnologias empregadas e os procedimentos experimentais com as respectivas métricas de avaliação.

% ---------------------------------------------------------------
\section{Tipo de pesquisa e hipótese}
\label{sec:tipo_pesquisa}

Informe qual tipo de pesquisa será conduzido (exploratória, empírica, formal etc.) e registre as hipóteses que serão avaliadas. Caso não haja hipótese formal, descreva o critério de sucesso esperado.

% ---------------------------------------------------------------
\section{Procedimentos da abordagem proposta}
\label{sec:procedimentos}

Detalhe as etapas da metodologia em ordem cronológica. Use listas numeradas, diagramas de fluxo ou pseudocódigo para tornar o processo claro e reprodutível.

% ---------------------------------------------------------------
\section{Ferramentas, tecnologias e ambientes}
\label{sec:ferramentas}

Liste softwares, linguagens, bibliotecas, frameworks e ambientes (laboratórios, nuvem, simulação) que serão usados. Justifique a escolha com base no problema e nas métricas definidas.

% ---------------------------------------------------------------
\section{Procedimentos experimentais e métricas}
\label{sec:experimentos_metricas}

Explique como os experimentos serão configurados, quais variáveis serão controladas e quais métricas demonstram que a proposta cumpre os objetivos (ex.: acurácia, latência, custo). Inclua tabelas que antecipem o plano de coleta, por exemplo:

\begin{table}[htb]
  \centering
  \caption{Plano de experimentação}
  \label{tab:plano_experimentos}
  \begin{tabular}{lcc}
    \toprule
    \textbf{Experimento} & \textbf{Variável controlada} & \textbf{Métrica principal} \\
    \midrule
    Comparação com baseline & algoritmo & acurácia \\
    Escalabilidade & carga & tempo médio \\
    \bottomrule
  \end{tabular}
  \legend{Fonte: Elaborado pelo autor (2025)}
\end{table}

\textit{Propor algo não basta: demonstre como será validado.}

% ---------------------------------------------------------------
% DESENVOLVIMENTO
% ---------------------------------------------------------------
\chapter{Desenvolvimento / Implementação / Experimentos}

\section{Modelagem e arquitetura}
Apresente diagramas (UML, blocos ou fluxogramas) que mostram como os módulos se relacionam, identificando componentes críticos e como eles atendem ao problema e à metodologia.

\section{Detalhes de implementação}
Explique escolhas tecnológicas, estruturas de dados e integrações. Inclua trechos de código relevantes ou pseudoexemplos se esclarecerem decisões de pesquisa.

\section{Experimentos e cenários de teste}
Descreva os cenários planejados, parâmetros testados e por que cada experimento está alinhado aos objetivos. Liste os conjuntos de dados, variações e número de repetições.

\section{Coleta e tratamento de dados}
Explique como os dados foram coletados (logs, sensores, APIs), armazenados e preparados (limpeza, normalização, anonimização). Registre ferramentas usadas no pré-processamento.

\textit{O desenvolvimento só é pesquisa se estiver claramente ligado ao problema e à metodologia definida.}

% ---------------------------------------------------------------
% RESULTADOS
% ---------------------------------------------------------------
\chapter{Resultados e Análise}

\section{Apresentação dos resultados}
Mostre resultados com tabelas, gráficos e figuras. Diferencie:
\begin{itemize}
  \item \textbf{Figura:} imagens, diagramas conceituais ou arquiteturais que ilustram estrutura ou fluxo.
  \item \textbf{Gráfico:} plot de dados (linhas, barras, boxplots) que compara métricas obtidas.
\end{itemize}

Exemplo de tabela:
\begin{table}[htb]
  \centering
  \caption{Comparação de métricas reportadas}
  \begin{tabular}{lcc}
    \toprule
    Método & Acurácia & Tempo médio (ms) \\
    \midrule
    Proposta & 92\% & 180 \\
    Baseline & 88\% & 220 \\
    \bottomrule
  \end{tabular}
  \legend{Tabela exemplificando como comparar métricas.}
\end{table}

Ao inserir um gráfico (por exemplo `assets/figures/desempenho.png`), use:
\begin{verbatim}
\begin{figure}[htb]
  \centering
  \includegraphics[width=0.8\textwidth]{assets/figures/desempenho.png}
  \caption{Gráfico de comparação das métricas de acurácia.}
  \label{fig:grafico-desempenho}
\end{figure}
\end{verbatim}

\section{Análise crítica e comparação}
Discuta as descobertas com base na revisão bibliográfica e na hipótese. Evidencie ganhos e perdas.

\section{Discussão}
Relacione os resultados aos objetivos e à hipótese, explicitando se o critério de sucesso foi atingido.

\textit{Resultados devem ser analisados, não apenas exibidos.}

% ---------------------------------------------------------------
% CONCLUSÃO
% ---------------------------------------------------------------
\chapter{Conclusão}

\section{Retomada do problema e dos objetivos}
Retome o problema e destaque o que foi atingido.

\section{Contribuições para a Computação}
Explique o que a pesquisa trouxe de novo.

\section{Limitações e trabalhos futuros}
Liste limites e sugira próximos passos coerentes.

% ---------------------------------------------------------------
% OBSERVAÇÕES DO AUTOR
% ---------------------------------------------------------------
\chapter*{Observações importantes}
\addcontentsline{toc}{chapter}{Observações importantes}

Para graduação, Wazlawick aceita trabalhos aplicados desde que:
\begin{itemize}
    \item O problema seja preciso;
    \item A solução esteja fundamentada;
    \item Haja análise crítica;
    \item Não seja apenas entrega funcional.
\end{itemize}

% ---------------------------------------------------------------
% APRESENTAÇÃO ORAL (BANCA)
% ---------------------------------------------------------------
\chapter*{Apresentação oral (banca)}
\addcontentsline{toc}{chapter}{Apresentação oral (banca)}

Além do documento escrito, o TCC2 é avaliado por uma \textbf{apresentação oral} final perante a banca. Um template pronto em \texttt{Beamer} está disponível no arquivo \texttt{apresentacao\_tcc2.tex}, na mesma pasta deste documento.

\noindent\fbox{\parbox{\dimexpr\textwidth-2\fboxsep-2\fboxrule}{%
\textbf{Regras da apresentação do TCC2}

\vspace{0.2cm}
\begin{itemize}
    \item \textbf{Duração: 30 minutos.} Cronometre os ensaios; ultrapassar o tempo é penalizado. A arguição da banca ($\approx$ 5 min) ocorre após a apresentação, separada desse total.
    \item \textbf{Cubra o trabalho completo}, seguindo a estrutura do documento escrito: introdução e problema, objetivos e método, desenvolvimento, \textbf{resultados e análise} e conclusão/contribuições.
    \item \textbf{Dê peso aos resultados:} grande parte do tempo deve ser dedicada a apresentar e \textbf{analisar criticamente} os resultados, não apenas descrevê-los.
    \item \textbf{Ritmo:} cerca de 1 a 1,5 minuto por slide, totalizando aproximadamente 20 a 26 slides de conteúdo.
    \item \textbf{Logo do SENAC:} grande no primeiro slide (capa) e pequena no canto superior direito dos demais slides --- já configurado no template.
    \item \textbf{Demonstração:} se houver demo ao vivo, tenha um plano B com capturas de tela ou vídeo curto.
    \item \textbf{Slides de apoio (backup):} inclua slides extras após o ``Obrigado'' com detalhes técnicos, trechos de código e tabelas completas para a arguição.
\end{itemize}

\vspace{0.2cm}
\textbf{Distribuição de tempo sugerida (30 min):} introdução/problema $\approx$ 4 min; objetivos/método $\approx$ 5 min; desenvolvimento $\approx$ 8 min; resultados/análise $\approx$ 8 min; conclusão $\approx$ 3 min; demonstração $\approx$ 2 min.
}}

% ---------------------------------------------------------------
% ELEMENTOS PÓS-TEXTUAIS
% ---------------------------------------------------------------
\postextual

% ===============================================================
% REFERÊNCIAS BIBLIOGRÁFICAS
% ===============================================================
% As referências bibliográficas são geradas automaticamente a partir
% do arquivo bibliografia_tcc2.bib, seguindo o padrão ABNT (NBR 6023).
%
% COMO CITAR no texto:
%   \cite{chave}         -> (SOBRENOME, ano) — citação indireta
%   \citeonline{chave}   -> Sobrenome (ano) — citação direta no texto
%
% TIPOS DE REFERÊNCIA no arquivo .bib:
%   @book        -> Livros
%   @article     -> Artigos de periódicos
%   @inproceedings -> Artigos em anais de eventos
%   @mastersthesis -> Dissertações de mestrado
%   @phdthesis   -> Teses de doutorado
%   @misc        -> Sites, vídeos e outros materiais
%   @manual      -> Documentação técnica
% ---------------------------------------------------------------
\bibliography{bibliografia_tcc2}

% ---------------------------------------------------------------
% APÊNDICES (se necessário)
% ---------------------------------------------------------------
% \begin{apendicesenv}
% \partapendices
%
% \chapter{Nome do Apêndice}
% \label{ap:apendice1}
%
% Conteúdo do apêndice (material elaborado pelo próprio autor).
%
% \end{apendicesenv}

% ---------------------------------------------------------------
% ANEXOS (se necessário)
% ---------------------------------------------------------------
% \begin{anexosenv}
% \partanexos
%
% \chapter{Nome do Anexo}
% \label{an:anexo1}
%
% Conteúdo do anexo (material de terceiros).
%
% \end{anexosenv}

\end{document}
